% Source: https://artofproblemsolving.com/community/c3h1221479_algebra_question_from_rice_exam_2014
% Stable identity: aops-topic-1221479-q01
A real cubic has exactly one real root (and two non-real ones) exactly when its discriminant is negative. For $x^3+bx^2+cx+d$ the discriminant is $18bcd-4b^3d+b^2c^2-4c^3-27d^2$; with $b=1$, $c=q$, $d=9$:
\[\Delta=162q-36+q^2-4q^3-2187=-(4q+39)\left(q^2-10q+57\right).\]
Since $q^2-10q+57=(q-5)^2+32>0$, $\Delta<0$ exactly when $q>-\dfrac{39}{4}$. At $q=-\dfrac{39}{4}$ there is a double root, so two distinct real roots. Hence $p=\boxed{-\dfrac{39}{4}}$.
